\section{Problem Setting}
\label{sec:setting}

We define a \textit{harness} $h$ as a persistent collection of tools, prompts, and skills that an agent can use to solve a task.
Given a task $t$ and a harness $h$, an agent can attempt the task using a loop of reasoning, acting, and observing.
This multi-step process generates a \textit{trajectory} $\tau$, which records the information read by the agent, its chain of thought, the tools used, and the final output.
We denote this execution process with a prompted agent operation as $\tau = \mathrm{solve}(h, t)$.
As the agent executes multiple tasks, it produces a dataset of trajectories $\mathcal{D} = \{\tau_1, \tau_2, \ldots, \tau_n\}$.
These trajectories often contain instances of failure and useful insights that can be used to improve the harness.
Consequently, we ask whether the agent can retrospectively analyze past trajectories to optimize its harness and improve its future performance.
To quantify this, we define a latent \textit{utility function} $U(t, \tau)$ that measures the quality of a trajectory.
We formalize the optimization as a function $\mathrm{optimize}(h, \text{instruction})$ that returns a modified harness $h'$.
The goal is to find an optimal harness $h^\star$ that maximizes the expected utility on future tasks:
\vspace{-0.5em}
\[
    h^\star = \arg\max_{h'} \; \mathbb{E}_{t,\, \tau \sim \mathrm{solve}(h', t)}
    \left[ U(t, \tau) \right].
\]

\noindent\textbf{Problem.}
However, estimating this utility function accurately is difficult in practice.
To evaluate the true utility of a harness, we would need a representative validation set of future tasks and a mechanism to calculate the success rate of the agent using this specific harness.
In our setting, the function $U$ is latent and cannot be directly observed.

\noindent\textbf{Our Approach.}
Because the utility $U$ is latent, we cannot directly optimize it.
Instead, we substitute this latent utility with a self-preference estimator.
Specifically, we instruct the agent to compare multiple trajectories on the same task to compute a self-preference ranking.
We define a ranking function as $(\text{rank}, \text{rationale}) = \mathrm{rank}(t, \tau_1, \tau_2, \ldots, \tau_m)$.
This function yields a preference ordering over the given trajectories and provides a rationale that explains why the agent prefers certain executions over others.
The next section details how we organize the operations of solving, ranking, and optimizing to improve the latent harness utility.